\paragraph{相位 Fourier 化的有限模态封口：Poisson--Cauchy 二阶形状函数的 Laurent 截断与矩反演}\label{para:cdim-poisson-cauchy-phase-fourierization}
本段把 \S\ref{subsec:unit-circle-phase-fourier} 的相位 Fourier 接口直接作用到 Poisson--Cauchy 粗粒化的二阶形状函数上：
通过 Cayley 相位变量 \(W=\iota(y)\) 将有理剖面提升为单位圆上的有限 Laurent 模态，从而把若干在正文中以三角积分核验的常数计算，统一压缩为有限维 Fourier 代数运算；
并由最高模态的不可混叠性给出中心矩的直接反演公式。

\begin{proposition}[Poisson--Cauchy 导数比的有限 Laurent 模态：相位变量下的严格截断]\label{prop:cdim-poisson-cauchy-finite-laurent-modes}
\leanverified{paper\_cdim\_poisson\_cauchy\_finite\_laurent\_modes}
设
\[
P_t(x):=\frac{1}{\pi}\frac{t}{x^2+t^2}\qquad(t>0),
\qquad
G(y):=\frac{1}{\pi(1+y^2)}.
\]
定义 Cayley 相位变量
\[
W(y):=\iota(y)=\frac{1+\mathrm{i}y}{1-\mathrm{i}y}\in\TT
\qquad(y\in\RR),
\]
并对整数 \(m\ge 0\) 定义归一化导数比
\[
Q_m(y):=t^m\,\frac{\partial_x^{\,m}P_t(ty)}{P_t(ty)}.
\]
则：
\begin{enumerate}
  \item \(Q_m(y)\) 与 \(t\) 无关；
  \item 存在唯一系数族 \(\{c_{m,k}\}_{|k|\le m}\subset\CC\) 使得
  \[
  \boxed{\;
  Q_m(y)=\sum_{k=-m}^{m}c_{m,k}\,W(y)^k\;}
  \qquad(y\in\RR),
  \]
  因而 \(Q_m\) 属于有限维 Laurent 模态空间 \(\mathrm{span}\{W^k:|k|\le m\}\)；
  \item 最高模态系数满足刚性闭式
  \[
  c_{m,m}=\frac{\mathrm{i}^{m}m!}{2^{m}},
  \qquad
  c_{m,-m}=\overline{c_{m,m}}=\frac{(-\mathrm{i})^{m}m!}{2^{m}}.
  \]
\end{enumerate}
特别地，令
\[
u_2(y):=\frac{3y^2-1}{(1+y^2)^2},\qquad
u_3(y):=\frac{4y(y^2-1)}{(1+y^2)^3},\qquad
u_4(y):=\frac{5y^4-10y^2+1}{(1+y^2)^4},
\]
则它们在 \(W\) 上具有显式有限 Laurent 展开
\begin{align}
u_2(y)
&=-\frac14\bigl(W^2+W+W^{-1}+W^{-2}\bigr),\label{eq:cdim-poisson-u2-laurent}\\
u_3(y)
&=\frac{\mathrm{i}}{8}\bigl(W^3+2W^2+W-W^{-1}-2W^{-2}-W^{-3}\bigr),\label{eq:cdim-poisson-u3-laurent}\\
u_4(y)
&=\frac1{16}\bigl(W^4+3W^3+3W^2+W+W^{-1}+3W^{-2}+3W^{-3}+W^{-4}\bigr).\label{eq:cdim-poisson-u4-laurent}
\end{align}
\end{proposition}
\begin{proof}
\textbf{(1)} 由恒等式 \(P_t(ty)=\frac{1}{\pi t}\frac{1}{1+y^2}\) 与 \(\partial_x=(1/t)\partial_y\) 立得。
\par
\textbf{(2)} 记 \(f(y):=(1+y^2)^{-1}\)，则 \(\partial_y^{\,m}f(y)=p_m(y)(1+y^2)^{-(m+1)}\)，其中 \(p_m\in\ZZ[y]\) 为次数至多 \(m\) 的多项式（对 \(m\) 归纳即可）。
故 \(Q_m(y)=(1+y^2)\partial_y^{\,m}f(y)=p_m(y)/(1+y^2)^m\)。
利用
\[
y=-\mathrm{i}\frac{W-1}{W+1},
\qquad
1+y^2=\frac{4W}{(W+1)^2}
\qquad\bigl(W=\iota(y)\bigr),
\]
可将 \(p_m(y)\) 写为 \((W+1)^{-m}\) 乘以次数 \(\le m\) 的多项式，而 \((1+y^2)^{-m}=(W+1)^{2m}/(4^m W^m)\)。
相乘后得到 \(Q_m(y)=W^{-m}\) 乘以次数 \(\le 2m\) 的多项式，因而为次数界 \(|k|\le m\) 的 Laurent 多项式。
唯一性由 \(\{W^k\}_{k\in\ZZ}\) 在单位圆上的线性无关性给出。
\par
\textbf{(3)} 为给出最高模态系数，取分式分解
\[
f(y)=\frac{1}{1+y^2}=\frac{1}{2\mathrm{i}}\Bigl((y-\mathrm{i})^{-1}-(y+\mathrm{i})^{-1}\Bigr),
\]
从而
\[
f^{(m)}(y)=\frac{(-1)^m m!}{2\mathrm{i}}\Bigl((y-\mathrm{i})^{-(m+1)}-(y+\mathrm{i})^{-(m+1)}\Bigr).
\]
由 \(Q_m=f^{(m)}/f=(1+y^2)f^{(m)}\) 与 \(1+y^2=(y-\mathrm{i})(y+\mathrm{i})\)，得到
\[
Q_m(y)=\frac{(-1)^m m!}{2\mathrm{i}}
\Bigl((y+\mathrm{i})(y-\mathrm{i})^{-m}-(y-\mathrm{i})(y+\mathrm{i})^{-m}\Bigr).
\]
由 \(W=\iota(y)=\frac{1+\mathrm{i}y}{1-\mathrm{i}y}=-\frac{y-\mathrm{i}}{y+\mathrm{i}}\) 可得
\(
y-\mathrm{i}=-\frac{2\mathrm{i}W}{W+1}
\)、
\(
y+\mathrm{i}=\frac{2\mathrm{i}}{W+1}
\)。
代入并化简可将 \(Q_m\) 写成 \((W+1)^{m-1}W^{-m}\) 与 \((W+1)^{m-1}W\) 的线性组合；其中 \((W+1)^{m-1}W\) 的最高项为 \(W^{m}\) 且系数为 \(1\)，于是
\(
c_{m,m}=\mathrm{i}^{m}m!/2^{m}
\)。
由 \(Q_m\) 的实值性（\(y\in\RR\)）得 \(c_{m,-m}=\overline{c_{m,m}}=(-\mathrm{i})^{m}m!/2^m\)。
\par
最后，\eqref{eq:cdim-poisson-u2-laurent}--\eqref{eq:cdim-poisson-u4-laurent} 由将 \(y=-\mathrm{i}\frac{W-1}{W+1}\) 与 \(1+y^2=\frac{4W}{(W+1)^2}\) 代入并化简得到。证毕。
\end{proof}

\begin{corollary}[有限代数积分原则：\texorpdfstring{$G(y)\,\dd y$}{G(y)dy} 下的 Laurent 模态平均为零阶系数]\label{cor:cdim-poisson-cauchy-laurent-integrals-rational}
沿用命题 \ref{prop:cdim-poisson-cauchy-finite-laurent-modes} 的记号。
对任意有限 Laurent 多项式 \(L(W)=\sum_{k=-K}^{K}a_k W^k\)（\(a_k\in\CC\)）有严格恒等式
\[
\boxed{\;
\int_{\RR} L\!\bigl(W(y)\bigr)\,G(y)\,\dd y=a_0\;}
\]
并因此，当 \(L\) 的系数属于 \(\QQ\oplus \mathrm{i}\QQ\) 时，上式取值亦属于 \(\QQ\oplus \mathrm{i}\QQ\)。
特别地，对任意多项式 \(P\in\CC[U_2,U_3,U_4]\)，积分
\(
\int_{\RR}G(y)\,P(u_2(y),u_3(y),u_4(y))\,\dd y
\)
可化为有限 Fourier 系数的代数运算，并落入 \(\QQ\oplus \mathrm{i}\QQ\)。
\end{corollary}
\begin{proof}
由 \S\ref{subsec:unit-circle-phase-fourier} 的变换 \(y=\tan(\pi u)\)（见 \eqref{eq:unit-circle-phase-fourier-coeff}），有 \(W(y)=e^{2\pi\mathrm{i}u}\) 且 \(G(y)\dd y=\dd u\)。
因此
\[
\int_{\RR} W(y)^k\,G(y)\,\dd y=\int_{-1/2}^{1/2}e^{2\pi\mathrm{i}ku}\,\dd u=\mathbf 1_{\{k=0\}},
\]
代入 \(L(W)=\sum a_k W^k\) 即得第一式。
对 \(P(u_2,u_3,u_4)\) 的结论由 \eqref{eq:cdim-poisson-u2-laurent}--\eqref{eq:cdim-poisson-u4-laurent} 使其化为 Laurent 多项式并应用上式得到。证毕。
\end{proof}

\begin{theorem}[\texorpdfstring{$L^p$}{Lp} 范数的锐渐近常数：二阶 Poisson--Cauchy 形状函数的全谱极限]\label{thm:cdim-poisson-cauchy-lp-sharp-constants}
令 \(\nu\) 为 \(\RR\) 上概率测度，重心 \(\bar\gamma\)，方差 \(\sigma^2\in(0,\infty)\)，并满足 \(\EE[|\gamma-\bar\gamma|^3]<\infty\)。
定义
\[
h_t(x):=\int_{\RR}P_t(x-\gamma)\,\dd\nu(\gamma),
\qquad
g_t(x):=P_t(x-\bar\gamma).
\]
则对任意 \(p\in(1,\infty)\) 存在严格极限
\[
\boxed{\;
\lim_{t\to\infty}\frac{t^{3-\frac1p}}{\sigma^2}\,\|h_t-g_t\|_{L^p(\RR)}
=
\frac1\pi\left(\int_{\RR}\frac{|3y^2-1|^p}{(1+y^2)^{3p}}\,\dd y\right)^{1/p}\;}
\]
并且端点范数满足
\[
\boxed{\;
\lim_{t\to\infty}\frac{t^{3}}{\sigma^2}\,\|h_t-g_t\|_{L^\infty(\RR)}=\frac1\pi\; }.
\]
特别地，
\[
\lim_{t\to\infty}\frac{t^{5/2}}{\sigma^2}\,\|h_t-g_t\|_{L^2(\RR)}=\frac{\sqrt3}{4\sqrt\pi}.
\]
\leanverified{paper\_cdim\_poisson\_cauchy\_lp\_sharp\_constants}
\end{theorem}
\begin{proof}
由命题 \ref{prop:cdim-poisson-l1-second-order-constant} 的二阶展开与
\(\|P_t^{(3)}\|_{L^1}=O(t^{-3})\)、\(\|P_t^{(3)}\|_{L^\infty}=O(t^{-4})\)，可得存在余项 \(\varepsilon_t\) 使
\[
h_t-g_t=\sigma^2 r_t+\varepsilon_t,
\qquad
r_t(x):=\frac12\,\partial_x^2 g_t(x),
\qquad
\|\varepsilon_t\|_{L^1}=O(t^{-3}),\ \ \|\varepsilon_t\|_{L^\infty}=O(t^{-4}).
\]
由插值不等式得对 \(p\in(1,\infty)\)，
\(
\|\varepsilon_t\|_{L^p}\le \|\varepsilon_t\|_{L^1}^{1/p}\|\varepsilon_t\|_{L^\infty}^{1-1/p}
=O(t^{-4+1/p})
=o(t^{-(3-1/p)}).
\)
另一方面，令 \(y=(x-\bar\gamma)/t\) 则
\[
r_t(\bar\gamma+ty)=\frac1{\pi t^3}\frac{3y^2-1}{(1+y^2)^3},
\qquad
\dd x=t\,\dd y,
\]
从而
\[
\|r_t\|_{L^p(\RR)}
=t^{-(3-1/p)}\cdot \frac1\pi\left(\int_{\RR}\frac{|3y^2-1|^p}{(1+y^2)^{3p}}\,\dd y\right)^{1/p}.
\]
合并并除以 \(\sigma^2\) 即得 \(p\in(1,\infty)\) 的极限常数。
\(\infty\)-范数由同一剖面式与 \(\sup_{y}\frac{|3y^2-1|}{(1+y^2)^3}=1\)（在 \(y=0\) 处取到）得到。
最后一式取 \(p=2\) 并用积分恒等式
\(
\int_{\RR}\frac{(3y^2-1)^2}{(1+y^2)^6}\dd y=\frac{3\pi}{16}
\)
即得。证毕。
\end{proof}

\begin{proposition}[相位 Fourier 系数的矩反演：由有限模态直接恢复 \texorpdfstring{$(\sigma^2,\mu_3,\mu_4)$}{(sigma^2,mu3,mu4)}]\label{prop:cdim-poisson-cauchy-moment-inversion-by-phase-modes}
在定理 \ref{thm:cdim-poisson-cauchy-lp-sharp-constants} 的设定下，进一步假设 \(\EE[|\gamma-\bar\gamma|^5]<\infty\)，并记中心矩
\[
\mu_3:=\EE[(\gamma-\bar\gamma)^3],\qquad \mu_4:=\EE[(\gamma-\bar\gamma)^4].
\]
令
\[
\delta_t(x):=\frac{h_t(x)}{g_t(x)}-1,
\qquad
y=\frac{x-\bar\gamma}{t},
\qquad
G(y)=\frac{1}{\pi(1+y^2)},
\qquad
W(y)=\iota(y).
\]
对 \(k\in\ZZ\) 定义相位 Fourier 系数
\[
\widehat{\delta_t}(k):=\int_{\RR}\delta_t(\bar\gamma+t y)\,\overline{W(y)^k}\,G(y)\,\dd y.
\]
则以下三个极限存在且为完全显式常数：
\[
\boxed{\;
\lim_{t\to\infty}t^2\,\widehat{\delta_t}(2)=-\frac{\sigma^2}{4},\qquad
\lim_{t\to\infty}t^3\,\widehat{\delta_t}(3)=\frac{\mathrm{i}}{8}\mu_3,\qquad
\lim_{t\to\infty}t^4\,\widehat{\delta_t}(4)=\frac{1}{16}\mu_4\;}
\]
从而可由三条模态的读数反演中心矩：
\[
\boxed{\;
\sigma^2=-4\lim_{t\to\infty}t^2\,\widehat{\delta_t}(2),\qquad
\mu_3=-8\mathrm{i}\lim_{t\to\infty}t^3\,\widehat{\delta_t}(3),\qquad
\mu_4=16\lim_{t\to\infty}t^4\,\widehat{\delta_t}(4)\;}
\]
（右端均为实数）。
\leanverified{paper\_cdim\_poisson\_cauchy\_moment\_inversion\_by\_phase\_modes}
\end{proposition}
\begin{proof}
将 \(h_t\) 以 \(\Delta:=\gamma-\bar\gamma\) 为变量在 \(\Delta=0\) 处作四阶 Taylor 展开并除以 \(g_t\)，可得（于 \(G(y)\dd y\) 的 \(L^3\) 意义下）
\[
\delta_t(\bar\gamma+t y)
=\frac{\sigma^2}{t^2}u_2(y)+\frac{\mu_3}{t^3}u_3(y)+\frac{\mu_4}{t^4}u_4(y)+O(t^{-5}),
\]
其中 \(u_2,u_3,u_4\) 如命题 \ref{prop:cdim-poisson-cauchy-finite-laurent-modes} 所定义。
由 \eqref{eq:cdim-poisson-u2-laurent}--\eqref{eq:cdim-poisson-u4-laurent} 知：
\(
u_2
\)
仅含 \(|k|\le 2\) 模态且其 \(W^2\) 系数为 \(-1/4\)；
\(
u_3
\)
仅含 \(|k|\le 3\) 模态且其 \(W^3\) 系数为 \(\mathrm{i}/8\)；
\(
u_4
\)
仅含 \(|k|\le 4\) 模态且其 \(W^4\) 系数为 \(1/16\)。
另一方面，\(\int_{\RR}W(y)^{j-k}G(y)\dd y=\mathbf 1_{\{j=k\}}\)（推论 \ref{cor:cdim-poisson-cauchy-laurent-integrals-rational} 的证明）给出 Fourier 系数的严格抽取。
故对 \(k=2,3,4\) 分别取 \(\widehat{\delta_t}(k)\) 并乘以 \(t^k\)，主项极限即为相应最高模态系数乘以中心矩，而其余项因阶数原因均为 \(o(1)\)。
整理即得三条极限与反演公式。证毕。
\end{proof}

\endinput
